\chapter{Source-Faithful PGVWC07 Recordings}\label{ch:pgvwc07_recordings}

This appendix-style chapter records the intermediate construction steps of the
single-tensor canonical-form theorem of
\cite[Theorem~Th:TIcanonical]{PerezGarcia2007Matrix}, following the source proof
order. They faithfully transcribe the PGVWC07 single-tensor construction and are
kept here only as a reference. None of these recordings is on the main route of
the Fundamental Theorem reduction, which instead goes through
Theorem~\ref{thm:tp_gauge_arbitrary} and the after-blocking statements in
Section~\ref{sec:reduction_to_normal_form}.

\section{Source-faithful PGVWC07 intermediate steps}\label{sec:pgvwc07_intermediates}

The theorems below record the intermediate construction steps of the proof of
\cite[Theorem~Th:TIcanonical]{PerezGarcia2007Matrix}, in the source proof order:
single-block unital and dual-diagonal construction; blockwise composition;
arbitrary-input form with the explicit zero summand; the length-zero dimension
identity; the positive-length form; the positive-length witness; weight
normalization; the nonzero-coefficient existence theorem; and the dichotomy that
yields Theorem~\ref{thm:pgvwc07_ti_canonical_form}. These are recordings of the
source structure. They are not used by the canonical-form reduction in the proof
of the Fundamental Theorem, which goes through
Theorem~\ref{thm:tp_gauge_arbitrary} and the after-blocking statements in
Section~\ref{sec:reduction_to_normal_form}.

\begin{theorem}[Single irreducible-block PGVWC07 canonical-form theorem]%
    \label{thm:pgvwc07_irreducible_unital_dual_package}
    \lean{MPSTensor.exists_pgvwc07_unital_dualDiag_data_of_irreducible}
    \leanok
    \uses{def:irreducible_tensor}
    Let $A$ be an irreducible nonzero MPS block with positive bond dimension.
    Then the Perron--Frobenius eigenvector produces a positive scalar $r$ and a
    positive definite matrix $\rho$ such that
    \[
        B^i := r^{-1/2}\rho^{-1/2} A^i\rho^{1/2}
    \]
    is unital. In this unital gauge every fixed point of $\E_B$ is a scalar
    multiple of the identity. Moreover a unitary conjugate $C^i=U^\dagger
        B^iU$ remains unital and has a diagonal positive definite dual fixed point
    $\Lambda$:
    \[
        \sum_i C^i(C^i)^\dagger=\Id,
        \qquad
        \sum_i (C^i)^\dagger\Lambda C^i=\Lambda.
    \]
    The construction preserves the finite-ring coefficients up to the
    spectral-radius weight carried by the rescaling from $A$ to $B$.

    This composes the single-block canonical-form existence steps of
    \cite[Theorem~Th:TIcanonical]{PerezGarcia2007Matrix}. It is not yet the
    full theorem, because the recursive finite-ring block decomposition and the
    final bond-dimension bound have not been combined into this construction.
\end{theorem}
\begin{proof}\leanok
    \uses{thm:unital_data_irreducible,
        thm:irreducible_unital_scalar_fixed_point,
        thm:pgvwc07_dual_diag_unital}
    Theorem~\ref{thm:unital_data_irreducible} gives the unital gauge.
    Gauge equivalence preserves irreducibility of the bond-space action, so
    Theorem~\ref{thm:irreducible_unital_scalar_fixed_point} shows that every
    fixed point of the unital representative is scalar. Then
    Theorem~\ref{thm:pgvwc07_dual_diag_unital} gives the unitary
    diagonalization of the dual fixed point.
\end{proof}

\begin{theorem}[Blockwise PGVWC07 unital and dual-diagonal block decomposition]%
    \label{thm:pgvwc07_blockwise_unital_dual_package}
    \lean{MPSTensor.exists_pgvwc07_unital_dualDiag_blockwise}
    \leanok
    \uses{def:irreducible_tensor,
        thm:pgvwc07_irreducible_unital_dual_package}
    Let a tensor~$A$ be represented, in finite-ring MPV coefficients, by a
    finite direct sum of nonzero irreducible blocks with unit weights. Then
    there is a weighted direct sum with the same MPV family such that every
    block is in the unital orientation,
    \[
        \sum_i C_k^i(C_k^i)^\dagger=\Id ,
    \]
    every fixed point of the block transfer map is a scalar multiple of the
    identity, and each block has a diagonal positive definite dual fixed point
    \[
        \sum_i (C_k^i)^\dagger\Lambda_k C_k^i=\Lambda_k.
    \]
    The weights are the positive spectral-radius weights produced by applying
    the single-block Perron--Frobenius gauge to each summand.

    This is the blockwise composition of
    \cite[Theorem~Th:TIcanonical]{PerezGarcia2007Matrix} after the recursive
    splitting has already produced the nonzero irreducible summands. It is not
    yet the full source theorem, because it does not start from an arbitrary
    representation and does not prove the final total bond-dimension bound.
\end{theorem}
\begin{proof}\leanok
    \uses{thm:pgvwc07_irreducible_unital_dual_package}
    Apply Theorem~\ref{thm:pgvwc07_irreducible_unital_dual_package} to each
    nonzero irreducible block. Gauge equivalence and the final unitary
    conjugation preserve the block MPV coefficients up to the displayed
    spectral-radius weights, and unitary conjugation carries the scalar
    fixed-point property of the unital transfer map.
\end{proof}

\begin{theorem}[Arbitrary-input PGVWC07 unital dual-diagonal form with zero block]%
    \label{thm:pgvwc07_arbitrary_zero_tail_unital_dual_package}
    \lean{MPSTensor.exists_pgvwc07_unital_dualDiag_from_arbitrary_with_zeroTail}
    \leanok
    \uses{thm:zero_block_separation,
        thm:pgvwc07_blockwise_unital_dual_package, def:same_mpv2_pos}
    Let $A$ be an arbitrary translation-invariant MPS tensor. Then there is
    a zero-block bond dimension $D_0$, a finite family of nonzero blocks
    $C_k$, and positive real weights $\mu_k$, such that for every chain
    length $N \ge 1$ and every word $\sigma \in \{0,\ldots,d-1\}^N$,
    \[
        \operatorname{MPV}_A
        =
        \operatorname{MPV}_{\oplus_k \mu_k C_k},
    \]
    together with the length-zero dimension identity $D = D_0 + \sum_k D_k$.
    Each nonzero block is in the unital orientation,
    \[
        \sum_i C_k^i(C_k^i)^\dagger=\Id ,
    \]
    every fixed point of its transfer map is a scalar multiple of the identity,
    and it has a diagonal positive definite dual fixed point
    \[
        \sum_i (C_k^i)^\dagger\Lambda_k C_k^i=\Lambda_k.
    \]
\end{theorem}
\begin{proof}\leanok
    \uses{thm:zero_block_separation,
        thm:pgvwc07_blockwise_unital_dual_package}
    First apply the zero-block separation theorem to $A$, retaining the
    all-zero summands only through their total bond dimension $D_0$ and
    extracting the nonzero irreducible summands. Then apply the blockwise
    PGVWC07 unital and dual-diagonal theorem to this nonzero family. Composing
    the positive-length equalities and the dimension identities gives the
    statement.
\end{proof}

\begin{theorem}[Bond-dimension bound for the zero-block PGVWC07 form]%
    \label{thm:pgvwc07_arbitrary_zero_tail_unital_dual_bond_dim_bound}
    \lean{MPSTensor.exists_pgvwc07_unital_dualDiag_from_arbitrary_with_zeroTail_bondDimBound}
    \leanok
    \uses{thm:pgvwc07_arbitrary_zero_tail_unital_dual_package}
    In the decomposition of
    Theorem~\ref{thm:pgvwc07_arbitrary_zero_tail_unital_dual_package}, the
    zero-block bond dimension and the nonzero block dimensions satisfy
    \[
        D = D_0+\sum_k D_k.
    \]
    In particular, the total bond dimension of the nonzero blocks is at most
    the original bond dimension $D$.
\end{theorem}
\begin{proof}\leanok
    \uses{thm:pgvwc07_arbitrary_zero_tail_unital_dual_package}
    The dimension identity is part of
    Theorem~\ref{thm:pgvwc07_arbitrary_zero_tail_unital_dual_package}, and the
    inequality follows because $D_0\geq 0$.
\end{proof}

\begin{theorem}[Positive-length PGVWC07 unital dual-diagonal form]%
    \label{thm:pgvwc07_arbitrary_positive_length_unital_dual_bond_dim_bound}
    \lean{MPSTensor.exists_pgvwc07_unital_dualDiag_from_arbitrary_posMPV_bondDimBound}
    \leanok
    \uses{thm:pgvwc07_arbitrary_zero_tail_unital_dual_bond_dim_bound}
    For any MPS tensor $A$ of bond dimension $D$, there are nonzero blocks
    $C_k$ and positive real weights $\mu_k$ such that each block is in the
    unital orientation,
    \[
        \sum_i C_k^i(C_k^i)^\dagger=\Id,
    \]
    has only scalar fixed points for its transfer map, and has a diagonal
    positive-definite fixed point for the dual transfer map,
    \[
        \sum_i (C_k^i)^\dagger\Lambda_k C_k^i=\Lambda_k.
    \]
    Moreover, for
    every chain length $N>0$ and word
    $\sigma\in\{0,\ldots,d-1\}^N$,
    \[
        \operatorname{MPV}_A(\sigma)
        =
        \operatorname{MPV}_{\bigoplus_k \mu_k C_k}(\sigma),
    \]
    and the nonzero block dimensions satisfy $\sum_k D_k\leq D$.
\end{theorem}
\begin{proof}\leanok
    \uses{thm:pgvwc07_arbitrary_zero_tail_unital_dual_bond_dim_bound}
    Apply Theorem~\ref{thm:pgvwc07_arbitrary_zero_tail_unital_dual_bond_dim_bound}.
    The all-zero summand contributes zero to every positive-length coefficient,
    so it may be omitted from the displayed MPV identity for $N>0$. The
    bound $\sum_k D_k\leq D$ is one of the conclusions of that theorem.
\end{proof}

\begin{definition}[Positive-length PGVWC07 canonical-form witness]%
    \label{def:pgvwc07_positive_length_witness}
    \lean{MPSTensor.PGVWC07PositiveLengthWitness}
    \leanok
    A positive-length PGVWC07 canonical-form witness for an MPS tensor $A$
    consists of a finite family of nonzero unital blocks, positive real
    weights, diagonal positive-definite dual fixed points, scalar fixed-point
    conclusions, positive block dimensions, positive-length MPV equality with
    $A$, and the total bond-dimension bound $\sum_k D_k\leq D$.

    This definition records the exact-MPV, nonempty-ring part of the PGVWC07
    construction. It does not include the source theorem's upper normalization
    $1\geq\mu_k$, which depends on the global spectral-radius normalization
    convention of \cite[Theorem~Th:TIcanonical]{PerezGarcia2007Matrix}.
\end{definition}

\begin{theorem}[Weight normalization for a nonempty PGVWC07 witness]%
    \label{thm:pgvwc07_positive_length_witness_weight_normalization}
    \lean{MPSTensor.PGVWC07PositiveLengthWitness.exists_weight_normalization}
    \leanok
    \uses{def:pgvwc07_positive_length_witness}
    Let a positive-length PGVWC07 canonical-form witness have at least one
    nonzero block. Then its positive real weights may be divided by their
    largest value so that the new weights $\nu_k$ are positive,
    $|\nu_k|\leq 1$ for all $k$, and $|\nu_{k_0}|=1$ for at least one
    block $k_0$. If the original weights are $\mu_k=s\nu_k$ with
    $s>0$, then for every chain length $N>0$,
    \[
        \operatorname{MPV}_A(\sigma)
        =
        s^N
        \operatorname{MPV}_{\bigoplus_k \nu_k C_k}(\sigma).
    \]
    This is the scalar normalization corresponding to
    \cite[Theorem~Th:TIcanonical]{PerezGarcia2007Matrix}; the displayed factor
    records the global length-dependent rescaling.
\end{theorem}
\begin{proof}\leanok
    \uses{def:pgvwc07_positive_length_witness,
        lem:mpv_to_tensor_from_blocks_weight_mul_left}
    Choose the largest positive real weight $s$. Dividing every weight by
    $s$ gives positive weights bounded above by $1$, and a block attaining
    the maximum has normalized weight $1$. The MPV identity follows from
    Lemma~\ref{lem:mpv_to_tensor_from_blocks_weight_mul_left}.
\end{proof}

\begin{theorem}[Projective form of PGVWC07 weight normalization]%
    \label{thm:pgvwc07_positive_length_witness_weight_normalization_projective}
    \lean{MPSTensor.PGVWC07PositiveLengthWitness.exists_weight_normalization_projective}
    \leanok
    \uses{def:pgvwc07_positive_length_witness,
        def:nonzero_proportional_mpv,
        thm:pgvwc07_positive_length_witness_weight_normalization}
    Let a positive-length PGVWC07 canonical-form witness have at least one
    block. After dividing the positive real weights by their largest value,
    the normalized weights $\nu_k$ are positive real numbers, satisfy
    $|\nu_k|\leq 1$ for all $k$, and satisfy $|\nu_{k_0}|=1$ for some block
    $k_0$. Moreover the original tensor and the normalized
    weighted block tensor generate nonzero proportional MPV families:
    for every length $N$ there is a nonzero scalar $c_N$ such that
    \[
        \operatorname{MPV}_A(\sigma)
        =
        c_N
        \operatorname{MPV}_{\bigoplus_k \nu_k C_k}(\sigma).
    \]
\end{theorem}
\begin{proof}\leanok
    \uses{thm:pgvwc07_positive_length_witness_weight_normalization,
        lem:mpv_zero_length}
    For $N>0$, take $c_N=s^N$, which is nonzero because $s>0$. For
    $N=0$, use the length-zero MPV coefficient and the positive total
    dimension of the nonzero blocks to choose the nonzero scalar relating the
    two traces.
\end{proof}

\begin{lemma}[Nonzero positive-length coefficient gives a nonempty PGVWC07 witness]%
    \label{lem:pgvwc07_positive_length_witness_nonempty_of_nonzero_mpv}
    \lean{MPSTensor.PGVWC07PositiveLengthWitness.block_count_pos_of_exists_ne_zero_mpv}
    \leanok
    \uses{def:pgvwc07_positive_length_witness,
        thm:mpv_decomposition}
    Let a positive-length PGVWC07 canonical-form witness for $A$ be given.
    If some positive-length MPV coefficient of $A$ is nonzero, then the
    witness has at least one block.
\end{lemma}
\begin{proof}\leanok
    \uses{def:pgvwc07_positive_length_witness,
        thm:mpv_decomposition}
    If the block family were empty, the block-diagonal MPV expansion would be
    an empty sum at every positive length. The positive-length MPV equality in
    the witness would then force all positive-length coefficients of $A$ to
    vanish, contradicting the hypothesis.
\end{proof}

\begin{theorem}[Nonzero projective PGVWC07 normalized form]%
    \label{thm:pgvwc07_nonzero_normalized_projective_form}
    \lean{MPSTensor.exists_pgvwc07_normalized_projective_form_of_exists_ne_zero_mpv}
    \leanok
    \uses{thm:pgvwc07_positive_length_witness,
        lem:pgvwc07_positive_length_witness_nonempty_of_nonzero_mpv,
        thm:pgvwc07_positive_length_witness_weight_normalization_projective}
    Suppose that some positive-length MPV coefficient of an MPS tensor $A$ is
    nonzero. Then $A$ has a nonempty finite family of nonzero PGVWC07 blocks
    $C_k$ with positive normalized weights $\nu_k$ satisfying
    $|\nu_k|\leq 1$ for all $k$ and $|\nu_{k_0}|=1$ for some block $k_0$.
    Each block is in the unital orientation, has only scalar fixed points for
    its transfer map, and has a diagonal positive-definite fixed point
    $\Lambda_k$ for the dual transfer map. Moreover the original tensor and
    the normalized weighted block tensor generate nonzero proportional MPV
    families, and the block dimensions satisfy $\sum_k D_k\leq D$.
\end{theorem}
\begin{proof}\leanok
    \uses{thm:pgvwc07_positive_length_witness,
        lem:pgvwc07_positive_length_witness_nonempty_of_nonzero_mpv,
        thm:pgvwc07_positive_length_witness_weight_normalization_projective}
    First use Theorem~\ref{thm:pgvwc07_positive_length_witness}. The nonzero
    positive-length coefficient makes the witness nonempty by
    Lemma~\ref{lem:pgvwc07_positive_length_witness_nonempty_of_nonzero_mpv}.
    Apply the projective weight-normalization theorem to this nonempty witness.
\end{proof}

\begin{theorem}[Exact normalized PGVWC07 form after global rescaling]%
    \label{thm:pgvwc07_nonzero_normalized_exact_rescaled_form}
    \lean{MPSTensor.exists_pgvwc07_normalized_exact_form_after_rescaling_of_exists_ne_zero_mpv}
    \leanok
    \uses{thm:pgvwc07_positive_length_witness,
        lem:pgvwc07_positive_length_witness_nonempty_of_nonzero_mpv,
        thm:pgvwc07_positive_length_witness_weight_normalization}
    Suppose that some positive-length MPV coefficient of an MPS tensor $A$ is
    nonzero. Then there is a positive scalar $s$ and a nonempty finite family
    of PGVWC07 blocks $C_k$ with positive normalized weights $\nu_k$
    satisfying $|\nu_k|\leq 1$ for all $k$ and
    $|\nu_{k_0}|=1$ for some block $k_0$. Each block is in the unital
    orientation, has only scalar fixed points for its transfer map, and has a
    diagonal positive-definite fixed point for the dual transfer map. Each
    block has positive bond dimension. Moreover the globally rescaled tensor
    $s^{-1}A$ and the normalized weighted block tensor have the same
    positive-length MPV coefficients, and the block dimensions satisfy
    $\sum_k D_k\leq D$.
\end{theorem}
\begin{proof}\leanok
    \uses{thm:pgvwc07_positive_length_witness,
        lem:pgvwc07_positive_length_witness_nonempty_of_nonzero_mpv,
        thm:pgvwc07_positive_length_witness_weight_normalization,
        lem:mpv_smul}
    Start from the positive-length witness and use the nonzero coefficient to
    make the block family nonempty. Divide the positive weights by their
    maximum. The exact coefficient identity for $A$ contains the global
    factor $s^N$ at length $N$; applying the scalar rescaling $s^{-1}$ to
    every matrix of $A$ contributes the factor $s^{-N}$, so the two factors
    cancel at every positive length.
\end{proof}

\begin{theorem}[Zero/nonzero dichotomy for exact rescaled PGVWC07 form]%
    \label{thm:pgvwc07_exact_rescaled_form_zero_nonzero_dichotomy}
    \lean{MPSTensor.exists_pgvwc07_normalized_exact_form_after_rescaling_or_forall_pos_mpv_eq_zero}
    \leanok
    \uses{thm:pgvwc07_nonzero_normalized_exact_rescaled_form}
    For every MPS tensor $A$, either every positive-length MPV coefficient of
    $A$ is zero, or the exact rescaled normalized PGVWC07 form of
    Theorem~\ref{thm:pgvwc07_nonzero_normalized_exact_rescaled_form} exists.
\end{theorem}
\begin{proof}\leanok
    \uses{thm:pgvwc07_nonzero_normalized_exact_rescaled_form}
    If some positive-length coefficient is nonzero, apply
    Theorem~\ref{thm:pgvwc07_nonzero_normalized_exact_rescaled_form}. Otherwise,
    every positive-length coefficient is zero.
\end{proof}

\begin{theorem}[Existence of the positive-length PGVWC07 witness]%
    \label{thm:pgvwc07_positive_length_witness}
    \lean{MPSTensor.exists_pgvwc07_positiveLengthWitness}
    \leanok
    \uses{def:pgvwc07_positive_length_witness,
        thm:pgvwc07_arbitrary_positive_length_unital_dual_bond_dim_bound}
    For every MPS tensor $A$, the conclusions of
    Theorem~\ref{thm:pgvwc07_arbitrary_positive_length_unital_dual_bond_dim_bound}
    give a positive-length PGVWC07 canonical-form witness.
\end{theorem}
\begin{proof}\leanok
    \uses{def:pgvwc07_positive_length_witness,
        thm:pgvwc07_arbitrary_positive_length_unital_dual_bond_dim_bound}
    Unpack
    Theorem~\ref{thm:pgvwc07_arbitrary_positive_length_unital_dual_bond_dim_bound}
    and use its conclusions as the fields of the witness.
\end{proof}
